\chapter{فلسفه‌ی بازنمایی}
\label{ch:philosophy-of-representation}

تا به حال به کلماتی مثل «طلا» یا «نقره» فکر کرده‌اید؟ یا چرا وقتی درختی را می‌بینیم، آن را «درخت» می‌نامیم، درحالی‌که خود درخت هیچ شباهتی به حروف و صدای این کلمه ندارد؟\\
چرا وقتی عدد «۵» می‌تواند هم به پنج شیء، هم به ۵ انسان، هم به پنج سال اشاره کند؟\\
چرا وقتی می‌گوییم «انقلاب فوریه» لازم نیست تمام جزئیات و رویدادهای آن را در ذهن بازسازی کنیم؟

وقتی به این‌ها فکر می‌کنیم، به نظر می‌رسد مغز و زبان ما راه عجیبی برای کنار آمدن با جهان یافته‌اند: به‌جای حفظ جزئیات دقیق اشیاء و رویدادها، با نماد‌ها و نام‌گذاری روی اشیاء کار می‌کنیم.

ما در انتهای فصل قبل سؤال کلیدی پرسیدیم: «چگونه یک کلمه، یک عدد یا یک نشانه می‌تواند این‌گونه جای چیز دیگری را بگیرد؟» هدف این فصل پاسخ به همین پرسش است. درواقع، این فصل با یکی از پایه‌ای‌ترین جنبه‌های طبیعت انسانی سروکار دارد: «نام‌گذاری» و نقش آن در به‌خاطر سپردن اطلاعات.

\section{چرا انسان به نام‌گذاری علاقه دارد؟}
\label{sec:why-naming}

نام‌گذاری و نمادگذاری سالیان سال همراه انسان بوده و هنوز هم او را همراهی می‌کند؛ از کشیدن اشکال مختلف روی دیوار غارها تا امروز که در کاغذ و فتوشاپ نقاشی می‌کشیم. اما چرا؟ چرا انسان این‌قدر علاقه دارد روی هرچیزی اسم بگذارد؟

ساده‌ترین جوابی که می‌توانیم به آن برسیم «تنبلی» است. ما باید قبول کنیم نژاد بشر بسیار تنبل‌تر از آن چیزی است که فکرش را می‌کند، و هرچقدر که ما تنبل هستیم، مغزمان تنبل‌تر است. برای روشن‌تر شدن ماجرا، فرض کنید مدیر یک انبار بزرگ هستید که روزانه ده‌ها هزار نامه‌ی طولانی دریافت می‌کنید و وظیفه دارید هرگاه کسی درخواست یک نامه را کرد، آن را در سریع‌ترین زمان ممکن پیدا و تحویل دهید. فکر می‌کنید بهترین راه چیست؟ پیش از خواندن ادامه‌ی متن، کمی تأمل کنید و بهترین راه را برای خودتان پیدا کنید و بعد به دنبال پاسخ بروید.

مغز ما همیشه با چنین وضعیتی روبه‌رو است. روزانه هزاران‌هزار چیز را باید در حافظه‌ی بلندمدت و کوتاه‌مدت ذخیره، مدیریت و بازیابی کند. اما چگونه؟

\subsection{مدیر انبار مغز: هیپوکامپ}
\label{sub:hippocampus-warehouse-manager}

یکی از اجزای مهم مغز برای مدیریت حافظه، بخش کوچکی به نام هیپوکامپ است که در انبار ما نقش مدیر انبار را دارد. نشریه‌ی کلیولند کلینیک در این باره توضیح می‌دهد که هیپوکامپ بخشی از مغز است که مسئول حافظه و یادگیری بوده، و این ساختار کوچک کمک می‌کند تا چه کوتاه‌مدت و چه بلندمدت، مطالب به خاطر سپرده شوند و از محیط آگاهی کسب شود\autocite{clevelandclinic-hippocampus}. به بیان دیگر، هیپوکامپ در تبدیل خاطرات کوتاه‌مدت به بلندمدت و سازمان‌دهی و بازیابی آن‌ها نقش مهمی دارد.

در کنار این «مدیر انبار»، شبکه‌ای کار می‌کند که مسئول بازیابی و مدیریت خاطرات است: حافظه‌ی اپیزودیک، یعنی توانایی رمزگذاری و بازیابی تجربیات شخصی روزانه، که توسط لوب گیجگاهی میانی (MTL) پشتیبانی می‌شود\autocite{roy2018hippocampalindex}. به تعبیر ساده، هیپوکامپ مانند مدیری عمل می‌کند که داده‌ها و نامه‌های حافظه را سروسامان می‌دهد، و شبکه‌ی حافظه‌ی اپیزودیک هم مانند سیستم پستی است که وظیفه‌ی جمع‌آوری و بازیابی خاطرات را دارد.

نکته‌ی مهم در سیستم بازیابی خاطره‌ی مغز این است که مغز تمام جزئیات هر رویداد را به شکل کامل حفظ نمی‌کند، بلکه بیشتر یک نمایه یا نماد کلی از رویداد را نگه می‌دارد. مثلاً اگر در تولدی کیک به صورتتان خورده باشد و باعث هیجان یا عصبانیت‌تان شده باشد، احتمالاً همین بخش به‌عنوان نماد کلی آن تولد در مغز شما ثبت شده است. به عبارت دیگر، حافظه‌ی ما شبیه به یک آرشیو کامل از صحنه‌ها نیست، بلکه مجموعه‌ای از ردِ حافظه‌هاست که به ما ارائه می‌شود\autocite{mitnews2017circuit}.

بااین‌حال، هدف ما ورود به جزئیات دقیق‌تر عملکرد مغز نیست؛ هدف، شناخت حالت کلی عملکرد ذخیره‌ی مغز بود، تا نشان دهد عملکرد حافظه برای زندگی روزمره حیاتی است و شامل انواع توانایی‌های خاص می‌شود که در هسته‌ی خود، امکان ذخیره و بازیابی اطلاعات را در بازه‌های زمانی متغیر، از ثانیه تا چند روز تا سال‌ها، فراهم می‌کند. علاوه‌براین، اکنون به‌خوبی شناخته‌شده است که اشکال متعددی از حافظه وجود دارد\autocite{frankland2019retrieval}.

و در عمل، وقتی بخواهیم خاطره‌ای مثل دوران «۹ سالگی» را به یاد بیاوریم، الزاماً کل اتفاقات آن دوره جلوی چشممان نمی‌آید، بلکه تنها چند نکته‌ی برجسته -- مثلاً یک ترس، هیجان یا لحظه‌ی خاص از دوران کودکی -- به خاطر می‌آید و ما به ذهن خود می‌گوییم: «آه، این مربوط به ۹ سالگی من بوده.» بر اساس این توضیحات درباره‌ی مغز، با چیزی آشنا می‌شویم به نام «بازنمایی» و «بازنمایی ذهنی».

\section{بازنمایی ذهنی}
\label{sec:mental-representation}

پیش از تعریف بازنمایی، می‌توانیم بازنمایی ذهنی را تعریف کنیم که دست‌کم تا ارسطو بازمی‌گردد. نقطه‌ی شروع خود را حالات ذهنی عقل سلیم، مانند افکار، باورها، آرزوها، ادراکات و تصورات، قرار می‌دهد. گفته می‌شود چنین حالاتی «قصدیت» دارند -- آن‌ها درباره‌ی چیزها هستند یا به آن‌ها اشاره می‌کنند و ممکن است با توجه به ویژگی‌هایی مانند سازگاری، حقیقت، تناسب و دقت ارزیابی شوند. برای مثال، این فکر که پسرعموها با هم نسبت فامیلی ندارند متناقض است، این باور که الویس مرده است درست است، میل به خوردن ماه نامناسب است، تجربه‌ی بصری توت‌فرنگی رسیده به رنگ قرمز دقیق است، و تصویرسازی جورج واشنگتن با موهای بافته نادرست است\autocite{sep-mental-representation}.

بازنمایی ذهنی همچنین فرایندهای ذهنی مانند تفکر، استدلال، و تخیل را به‌عنوان توالی حالات ذهنی عمدی درک می‌کند. برای مثال، تصور طلوع ماه بر فراز کوه، از جمله، به معنای در نظر گرفتن مجموعه‌ای از تصاویر ذهنی از ماه و یک کوه است\autocite{sep-mental-representation}.

از این منظر، درمی‌یابیم که حتی در بازیابی خاطرات هم با بازنمایی‌های ذهنی سروکار داریم، نه با خودِ گذشته. مثلاً وقتی می‌گوییم «درخت»، آیا واقعاً داریم خودِ درخت را می‌بینیم یا تصویر ذهنی‌ای که این کلمه در ذهن ما ساخته است؟ یا وقتی از «انقلاب فوریه» حرف می‌زنیم، آیا تمام رویدادها و تصاویر تاریخی در ذهن مرور می‌شوند یا صرفاً نامی است که برای بازنمایی آن مجموعه‌رویدادها انتخاب کرده‌ایم؟

سؤالات زیادی می‌توان پرسید، و به یک چیز عجیب رسیدیم: ما برای مدیریت جهان خودمان، الزاماً خودِ جهان را با تمام جزئیاتش با خود حمل نمی‌کنیم. برای درخت یک نام داریم، برای تعداد یک عدد، برای یک رویداد تاریخی یک عبارت، و حتی برای خاطرات خودمان هم چیزی که به ذهن برمی‌گردد الزاماً خودِ اتفاق نیست.

و وقتی دقت می‌کنیم، همه‌ی این‌ها یک چیز را به‌جای چیز دیگری قرار می‌دهند. این دقیقاً همان چیزی است که از ابتدا درباره‌اش حرف می‌زنیم: «بازنمایی»، که بازنمایی ذهنی مثال بسیار دقیقی از آن است.

\section{چهار معیار بازنمایی}
\label{sec:four-criteria-representation}

اما تا اینجا که پیش رفتیم، سؤالی بی‌پاسخ مانده است: وقتی می‌گوییم چیزی نماینده‌ی چیز دیگری است، دقیقاً چه چیزی باعث می‌شود آن را یک «بازنمایی» بدانیم؟ آیا هر چیزی که به چیز دیگری اشاره کند، بازنمایی محسوب می‌شود؟ مثلاً یک تکه کاغذ روی میز می‌تواند یادآور یک قرار باشد، اما آیا خودش یک بازنمایی است؟

فلسفه و علوم شناختی سال‌هاست با همین سؤال درگیرند. یکی از تلاش‌ها برای روشن‌تر کردن این مسئله، به‌خصوص در زمینه‌ی علوم اعصاب، چهار معیار برای مفهوم بازنمایی پیشنهاد می‌کند\autocite{baker2021representation}. این چهار معیار قرار نیست تعریف نهایی و مورد توافق همه‌ی فیلسوفان باشند، بلکه می‌توانند کمک کنند بفهمیم وقتی درباره‌ی یک بازنمایی صحبت می‌کنیم، دقیقاً چه چیزی مدنظرمان است.

\subsection{ارتباط با جهان واقعی}
\label{sub:criterion-real-world-connection}

نخست، بازنمایی باید با چیزی که نمایندگی می‌کند در جهان واقعی ارتباط داشته باشد. به این معنی که باید بتوانیم بین «نماینده» و «چیزی که نمایندگی شده» نوعی رابطه پیدا کنیم. اگر هیچ رابطه‌ای میان آن دو وجود نداشته باشد، سخت می‌شود گفت یکی بازنمای دیگری است.

\subsection{نقش علّی در رفتار}
\label{sub:criterion-causal-role}

دوم، بازنمایی باید بتواند در رفتار یا یک فرایند نقش علّی داشته باشد. یعنی بازنمایی نباید صرفاً همراه یا همبسته با چیزی باشد؛ باید بتواند در توضیح اینکه یک موجود یا یک سیستم چه کاری انجام می‌دهد نقشی داشته باشد. مثلاً اگر مغز چیزی را به‌عنوان «خطر» بازنمایی می‌کند، این بازنمایی باید در توضیح رفتاری مثل فرارکردن نقشی داشته باشد.

\subsection{انحصار نقش علّی}
\label{sub:criterion-exclusivity}

سوم، یک بازنمایی باید منحصراً نقش علّی خود را داشته باشد -- عدم وجود بازنمایی باید رفتاری را که قادر به انجام آن است مسدود کند. یعنی اگر چند بخش یا چند فرایند مختلف بتوانند دقیقاً همان نقش را در ایجاد یک رفتار ایفا کنند، دشوارتر می‌شود بگوییم یکی از آن‌ها به‌طور مشخص «بازنمایی» موردنظر ماست. به همین دلیل، در علوم اعصاب صرف اینکه یک الگوی فعالیت با یک ویژگی از محیط همبستگی داشته باشد، لزوماً کافی نیست؛ باید نقش مشخص آن الگو در سازوکار ایجاد رفتار هم روشن باشد.

\subsection{توجیه غایت‌شناختی}
\label{sub:criterion-teleological}

چهارم، یک بازنمایی باید توجیه غایت‌شناختی داشته باشد که مواردی را که در آن یک بازنمایی نادرست است مشخص کند. این معیار می‌گوید اگر چیزی واقعاً بازنمایی باشد، باید بتوانیم توضیح دهیم این بازنمایی قرار است چه کاری انجام دهد و چه زمانی می‌توانیم بگوییم این بازنمایی «اشتباه» کرده است. مثلاً یک نقشه می‌تواند مسیر را درست یا غلط نشان دهد؛ یک جمله می‌تواند درست یا نادرست باشد. پس برای اینکه چیزی را بازنمایی بدانیم، باید معیاری داشته باشیم که به ما اجازه دهد بین یک بازنمایی موفق و یک بازنمایی نادرست تفاوت بگذاریم.

\section{چیزی به‌جای چیز دیگر}
\label{sec:something-in-place-of-another}

در این فصل دیدیم که مغز و زبان ما چگونه از بازنمایی برای کار با جهان استفاده می‌کنند. به‌جای اینکه هر چیز را با تمام جزئیاتش نگه داریم، چیزی که «به‌جای» آن چیز است می‌سازیم. برای مثال، یک کلمه یا عدد خودِ شیء نیست، بلکه نماینده‌ای از آن شیء است. همین‌طور یک نقشه یا عکس، خودِ جهان نیستند، بلکه جایگاه آن را نشان می‌دهند.

احتمالاً این را هم در کودکی از مدیر مدرسه یا معلمی شنیده‌اید: «شما آینه‌ی پدر و مادرتان هستید.» خب، شما واقعاً آینه‌ی پدر و مادرتان نیستید، اما ممکن است رفتارها، عادت‌ها، حرف‌زدن یا حتی بعضی از ویژگی‌های آن‌ها را در خودتان بازتاب دهید. ما حتی در زبان روزمره هم دائماً از چیزی استفاده می‌کنیم که خودش چیز دیگری نیست، اما به‌نحوی ما را به آن چیز مربوط می‌کند.

درواقع، چیزی که در تمام این مثال‌ها مشترک است همین است: یک چیز در جای چیز دیگری قرار می‌گیرد. حال تا اینجا که فهمیدیم، بگذارید یک سؤال به‌عنوان تمرین به شما بدهم: آیا بازنمایی فقط یک مفهوم انتزاعی در ذهن ماست، یا می‌تواند در یک سیستم فیزیکی هم وجود داشته باشد؟ و اگر قرار است در جهان واقعی وجود داشته باشد، چگونه از یک مفهوم انتزاعی به چیزی فیزیکی و قابل اندازه‌گیری تبدیل می‌شود؟